\documentclass[a4paper]{article}

\usepackage{graphicx}
\usepackage{amsmath,amssymb}
\usepackage{minted}
\usepackage{multirow}
\usepackage{float}
\usepackage{todonotes}
\usepackage{forest}
\usepackage{xstring}
\usepackage{xcolor}
\usepackage[most]{tcolorbox}
\usepackage{xparse}
\usepackage{natbib}

\usetikzlibrary{graphs,graphdrawing}
\usegdlibrary{layered}

% for external graphviz
\input{/home/donkarlo/Dropbox/repo/nd_language_ai_project/src/nd_language_ai/formal/computer/programming/tex/latex/code_clips/external_graphviz.tex}

% for tags
\input{/home/donkarlo/Dropbox/repo/nd_language_ai_project/src/nd_language_ai/formal/computer/programming/tex/latex/parts/control_sequence/kind/semtag/semtag.tex}

% for links
\input{/home/donkarlo/Dropbox/repo/nd_language_ai_project/src/nd_language_ai/formal/computer/programming/tex/latex/code_clips/seehere.tex}

\title{Neural encoding}
\date{}
\author{Mohammad Rahmani}

\begin{document}
    \maketitle
    Neural encoding: from stimulus -> to neural response
    
    Neural encoding concerns the relationship between an external stimulus and the neural activity(response) that it evokes. The stimulus is systematically varied while responses from relevant neurons or neuronal populations are recorded. The resulting activity is analyzed to determine which spatial, temporal, or firing-rate features reliably represent stimulus properties. A candidate neural code must preserve enough information to distinguish between stimuli that the organism itself can discriminate. Thus, encoding identifies which measurable components of neural activity carry information about changes in the sensory input \seeherefooter{johnson-2000-neural-coding}.
    \bibliographystyle{plainnat}
    \bibliography{/home/donkarlo/Dropbox/repo/nd_shared_research/bibliography/bibliography.bib}
\end{document}